\thispagestyle{empty}

\vspace{-2cm}
\begin{center}
\textbf{\large \textcolor{black}{چکیده} }
\end{center}

چكيده بخشي از پايان‌نامه است كه خواننده را به مطالعه آن علاقمند مي‌كند و يا از آن مي‌گريزاند. چكيده بايد ترجيحاً‌ در يك صفحه باشد ( تقريباً تمامي چكيده پايان‌نامه‌ها در يك صفحه قابل نگارش است). در نگارش چكيده نكات زير بايد رعايت شود. متن چكيده بايد مزين به كلمه‌ها و عبارات سليس، آشنا، با معني و روشن باشد. بگونه ای که با  حدود 300 تا 500 کلمه بتواند خواننده را به خواندن پایان نامه راغب نماید. چكيده، جدا از پايان‌نامه بايد به تنهايي گويا و مستقل باشد. در چكيده باید از ذكر منابع، اشاره به جداول و نمودارها اجتناب شود. تميز بودن مطلب، نداشتن غلط‌هاي املايي يا دستور زباني و رعايت دقت و تسلسل روند نگارش چكيده از نكات مهم ديگري است كه بايد درنظر گرفته شود. در چكيده پايان‌نامه بايد از درج مشخصات مربوط به پايان‌نامه خودداري شود. کلمات كليدي در انتهاي چكيده فارسی و انگلیسی آورده شود. محتواي چكيده‌ها بر اساس موضوع و گرايش تحقيق طبقه‌بندي مي‌شود و به همين جهت وجود كلمات شاخص و كليدي، مراكز اطلاعاتي  را در طبقه‌بندي دقيق و سريع پایان‌نامه ياري مي‌دهد.  كلمات كليدي، راهنماي نكات مهم موجود در پايان‌نامه هستند. بنابراين بايد در حد امكان كلمه‌ها يا عباراتي انتخاب شد كه ماهيت، محتوا و گرايش كار را به وضوح روشن نمايد. چكيده بايد منعكس كننده اصل موضوع باشد. در چكيده بايد اهداف تحقیق مورد توجه قرار گيرد. تأكيد روي اطلاعات تازه ( يافته‌ها) و اصطلاحات جديد يا نظريه‌ها، فرضيه‌ها، نتايج و پيشنهادها متمركز شود. اگر در پايان‌نامه روش نويني براي اولين بار ارائه مي‌شود و تا به حال معمول نبوده است، با جزئيات بيشتري ذكر  شود. شايان ذكر است چكيده فارسی و انگلیسی بايد حتماً به تأييد استاد راهنما رسيده باشد.

\hspace{-0.65cm}%
\textbf{\textcolor{black}{کلمات کلیدی}} :
تعداد كلمات يا عبارات كليدي حداكثر مي‌تواند پنج كلمه يا عبارت باشد. کلید واژه‌ها باید با واژه‌های اصلی عنوان و مسئله تحقیق تناسب داشته باشند

\clearpage
